% Generated draft: requires mathematical and source review.
\documentclass[11pt]{article}
\usepackage[margin=25mm]{geometry}
\usepackage{fontspec}
\setmainfont{Latin Modern Roman}
\usepackage{amsmath,amssymb,mathtools,unicode-math}
\setmathfont{Latin Modern Math}
\usepackage{enumitem,needspace,xurl,hyperref}
\setlength{\parindent}{0pt}
\setlength{\parskip}{5pt}
\setlength{\emergencystretch}{3em}
\newcommand{\N}{\mathbb{N}}
\newcommand{\Z}{\mathbb{Z}}
\newcommand{\Q}{\mathbb{Q}}
\newcommand{\R}{\mathbb{R}}
\newcommand{\C}{\mathbb{C}}
\newcommand{\F}{\mathbb{F}}
\newcommand{\A}{\mathbb{A}}
\newcommand{\PP}{\mathbb{P}}
\newcommand{\E}{\mathbb{E}}
\newcommand{\eps}{\varepsilon}
\newcommand{\dd}{\,\mathrm{d}}
\begin{document}
\section*{UnsolvedMath ID: 20001127}
\textbf{Source problem code:} catalogue entry
\subsection{Definitions and mathematical statement}
Brought by T. Tao and originally stated by I. Ruzsa, the source asks whether $A\subseteq\F_2^n$ with $|A+A|\leq K|A|$ has a linear subspace $V$ satisfying
\[
 |V|<K^c|A|,\qquad |A\cap V|\geq K^{-c}|A|,
\]
for a fixed exponent. The strict sign and pure powers of $K$ are retained. The aid distinguishes this literal normalization from the solved polynomial Freiman--Ruzsa (PFR) statement.

At $K=1$, a non-linear affine coset $A=a+H$, $a\notin H$, is an immediate obstruction: the intersection requirement forces $A\subseteq V$, hence $H+\langle a\rangle\subseteq V$ and $|V|\geq2|A|$. Even a linear $A$ shows the strict size inequality alone cannot hold at $K=1$ together with full intersection.

The aid gives a robust family with $K>1$ tending to one. Let $|H|=q\geq4$ be a power of two, choose $a\notin H$ and set
\[
 A=a+(H\setminus\{0\}),\qquad |A|=q-1,\qquad A+A=H,
 \qquad K=\frac q{q-1}.
\]
For every fixed $c>0$, sufficiently large $q$ makes
\[
 K^{-c}(q-1)>q/2,\qquad K^c(q-1)<2q.
\]
A linear subspace meeting $A$ in more than $q/2$ points must contain both $H$ and $a$, so its size is at least $2q$, violating the other bound. Hence no fixed exponent repairs the literal two inequalities uniformly down to $K=1$.

The intended PFR problem is nevertheless reported solved. The summary invokes Liao's 2024 cover by at most $2K^9$ cosets of a subspace $H$ with $|H|\leq|A|$. For $K\geq2$, a coset meeting $A$ in at least $|A|/(2K^9)$ points is contained in its linear span $V=H+\langle a\rangle$, and the aid obtains the exact source inequalities with the single exponent $c=10$:
\[
 |V|<K^{10}|A|,\qquad |A\cap V|\geq K^{-10}|A|.
\]
The restriction $K\geq2$ absorbs the fixed factors and is essential to this stated corollary.

The original remark calls the following an equivalent formulation. For $f:\F_2^m\to\F_2^\infty$ let
\[
 S=\{f(x+y)-f(x)-f(y):x,y\in\F_2^m\}.
\]
Can one write $f=g+h$, with $g$ linear and $|\operatorname{im}h|$ polynomial in $|S|$? The aid says Gowers--Green--Manners--Tao, Corollary 1.4, proves this approximate-homomorphism statement. No explicit exponent is supplied. Equivalence here is with the properly normalized intended PFR theorem, not with the refuted all-$K$ strict inequalities word for word. The target $\F_2^\infty$ is understood as the direct-sum vector space; the finite domain gives finite image and defect set.
\subsection{Short English statement}
The literal linear-subspace bounds fail near doubling one because affine shifts require a fixed size factor. The aid separates that defect from the solved PFR theorem: a named covering bound gives the exact requested inequalities for $K\geq2$, and a named corollary proves the approximate-homomorphism formulation.
\subsection{Sources}
\begin{itemize}
\item Local UnsolvedMath registry, permanent ID $20001127$, source code \texttt{catalogue entry}, supplied complete source record, original remark and research summary.
\item American Institute of Mathematics, \emph{Recent trends in additive combinatorics}, Problem 3.15, brought by T. Tao and originally stated by I. Ruzsa; supplied URL \url{https://aimath.org/WWN/additivecomb/additivecomb.pdf}. Canonical local record \texttt{aim-combinatorics-notes.json}, \texttt{notes[251]}. The inserted page-10 footer credits collection to Ernie Croot and Vsevolod F. Lev; it is not part of the function's codomain or formula. No PDF was retrieved.
\item Liao's 2024 exponent-nine covering theorem and Gowers--Green--Manners--Tao, Corollary 1.4, as cited in the supplied summary. Full titles, identifiers and bibliographic entries are absent; the original literature field is empty.
\item Research attempt 1: counterexample to literal normalization; problem refuted at the run, registry status solved and classification \texttt{SOLVED-IN-LITERATURE}; supplied rigor label ``proved'', novelty confidence ``low''. The machine-generated aid is not peer reviewed and requires independent verification before citation.
\end{itemize}
\end{document}
